% EVAL FIXTURE — deliberately flawed. Invented content, not a template.
% Failure under test: stage 2. There is a topic but no nameable question, and no
% instability. Stages 3-7 are seeded with bait (fillers, passive with a human agent,
% and/then chains, an uncited claim) so a passing review must NOT report them.
\documentclass[11pt,a4paper]{article}
\usepackage[b1]{ercformatting}
\usepackage{ercreview}

\title{Coastal Sediment Transport Across Scales}
\acronym{SEDISCALE}
\author[Aro]{L. Aro}
\institution{Department of Earth Sciences, Example University, Country}

\begin{document}
\maketitle

\begin{abstract}
Coastal sediment transport is a very important topic in Earth science. This project
will study sediment transport at several scales and will develop a modelling
framework. The results will clearly be of interest to a wide range of communities.
\end{abstract}

\startsectiona

\section{Sediment transport across scales}
Coastal sediment transport has been studied for decades, and it is widely
believed to be a central control on shoreline evolution. Various approaches have been
applied, and the literature is extensive. Our group has worked on this topic in
several settings, and we have accumulated a large body of field observations.
Understanding sediment transport is essential for understanding coasts.

The topic spans many scales, from grain to shoreline. Each scale has its own
literature and its own instruments. A great deal remains to be learned.

\section{Vision and aim}
\textbf{Vision.} A quantitative understanding of how coastal landforms organise
themselves, independent of the particular coast being studied.

\textbf{Aim.} Within five years, to establish which grain-scale processes control
shoreline-scale organisation in wave-dominated coasts.

\section{Objectives}
\textbf{O1.} Identify which grain-scale transport processes propagate to
shoreline-scale morphology, and which are averaged out.

\textbf{O2.} Determine whether the scale coupling is set by sediment supply or by
wave climate, using contrasting field sites.

\textbf{O3.} Test whether shoreline-scale organisation can be predicted from
grain-scale measurements alone.

\section{State of the art}
Existing models operate at a single scale. Cross-scale coupling has been proposed but
never measured directly. Measurements will be made by us at both scales
simultaneously, and then the coupling will be quantified.

\section{Impact}
The findings will simply be relevant to geomorphology, coastal engineering, and
sedimentology.

\startsectionb

\cvsection*{Personal Details}
ORCID: 0000-0002-0000-0000

\cvsubsection{Current position(s)}
2021 -- present, Senior Lecturer, Example University

\cvsubsection{Ten selected outputs}
1. Aro L. et al. (2024). Grain-scale transport under breaking waves. \emph{J. Example
Geophys.} First author; first simultaneous measurement of grain flux and bar
migration; supplies the measurement principle for O1.

\end{document}
